% ======================================================================
%  PART 6 -- Fourier Theory, Spectral Density and the Periodogram
% ======================================================================
\section{Fourier Theory, Spectral Density and the Periodogram}
\label{sec:spectral}

The frequency domain is the second face of a stationary process. Every
second-order property carried by the autocovariance sequence (ACVS)
$\{\acvs_\tau\}$ in the time domain is mirrored, via the Fourier transform, by
the \textbf{spectral density function} (SDF) $\sdf(f)$ in the frequency domain.
This part builds that bridge in three stages. First, pure Fourier theory:
continuous-time, discrete-time and finite-data transforms, the convolution
theorem, aliasing and the Nyquist frequency. Second, the spectral
representation of a stationary process: $\sdf(f)=\sum_\tau\acvs_\tau
e^{-2\pi i f\tau}$, Herglotz's theorem, the integrated spectrum and its
Lebesgue decomposition, and the spectral representation theorem with its
orthogonal-increment process $\{Z(f)\}$. Third, estimation: the periodogram
$\hat\sdf^{(p)}(f)=\abs{J(f)}^2$, why it is (asymptotically) unbiased but
\emph{inconsistent}, and how tapering cures bias while multitapering cures
variance. Throughout we take the sampling interval $\Delta=1$ once we reach
the spectral chapters, and we always work with ordinary frequency $f$ (not
angular frequency $\omega=2\pi f$).

% ======================================================================
\subsection{Definitions}
% ======================================================================

\begin{definition}{Continuous-time Fourier transform and its inverse}{p6-ctft}
For $g:\Reals\to\Comp$ with $g\in L^1$, the \textbf{Fourier transform}
$G$ (upper case denotes the transform; some authors write $\hat g$) is
\[
G(f)=\int_{-\infty}^{\infty} g(t)\,e^{-2\pi i f t}\,dt,\qquad f\in\Reals .
\]
If in addition $G\in L^1$, the \textbf{inverse transform} recovers $g$ for
almost all $t$:
\[
g(t)=\int_{-\infty}^{\infty} G(f)\,e^{2\pi i f t}\,df .
\]
We write $g\longleftrightarrow G$ and call $(g,G)$ a \textbf{Fourier pair}.
\end{definition}

\begin{definition}{Properties of the continuous-time transform}{p6-ctprops}
For $g,h\in L^1$ with transforms $G,H$ and scalar $\alpha$, the basic pairs are
\[
\begin{aligned}
\conj{g(t)} &\longleftrightarrow \conj{G(-f)}, &\qquad
g(\alpha t)\;(\alpha\ne0) &\longleftrightarrow \tfrac{1}{\abs\alpha}G(f/\alpha),\\
g(t+\tau) &\longleftrightarrow G(f)\,e^{2\pi i \tau f}, &\qquad
\alpha g(t)+h(t) &\longleftrightarrow \alpha G(f)+H(f).
\end{aligned}
\]
(Conjugation reflects, scaling in time inversely scales frequency, a time
shift multiplies by a phase, and the transform is linear.)
\end{definition}

\begin{definition}{Continuous-time convolution}{p6-ctconv}
For $g,h\in L^1$, the \textbf{convolution} $u=g*h$ is
\[
u(t)=\int_{-\infty}^{\infty} g(t-u)\,h(u)\,du,\qquad t\in\Reals .
\]
Convolution at a point is a weighted average of one function by the other,
centred at that point.
\end{definition}

\begin{definition}{Discrete-time Fourier transform and its inverse}{p6-dtft}
For a sequence $\{g_t\}_{t\in\Delta\Ints}$ in $\ell^1$, the
\textbf{discrete-time Fourier transform} is
\[
G(f)=\Delta\sum_{t\in\Delta\Ints} g_t\,e^{-2\pi i t f},\qquad f\in\Reals,
\]
with inverse (a Fourier series with the roles of domains swapped)
\[
g_t=\int_{-1/(2\Delta)}^{1/(2\Delta)} G(f)\,e^{2\pi i t f}\,df,\qquad t\in\Delta\Ints .
\]
A defining feature is \textbf{periodicity}: $G(f)=G(f+k/\Delta)$ for every
$k\in\Ints$, so $G$ is determined by its values on one period of length
$1/\Delta$.
\end{definition}

\begin{definition}{Discrete-time and periodic convolution}{p6-dtconv}
For $g,h\in\ell^1$ the \textbf{discrete convolution} $u=g*h$ is
\[
u_t=\Delta\sum_{u\in\Delta\Ints} g_{t-u}\,h_u,\qquad t\in\Delta\Ints,
\]
(with $\Delta=1$ the prefactor disappears). The dual object, the
\textbf{periodic convolution} $U=G*H$ of two periodic spectra, is
\[
U(f)=\int_{-1/(2\Delta)}^{1/(2\Delta)} G(f-f')\,H(f')\,df' .
\]
\end{definition}

\begin{definition}{Finite-data Fourier transform and the Dirichlet kernel}{p6-fdft}
Given observations $g_0,\dots,g_{(n-1)\Delta}$ at times
$T=\{0,\Delta,\dots,(n-1)\Delta\}$, the \textbf{finite Fourier transform} is
\[
G_n(f)=\sum_{t\in T} g_t\,e^{-2\pi i t f},\qquad f\in\Reals .
\]
It equals the full transform of $g_t$ multiplied by the indicator
$d_t=\tfrac1\Delta\ind_T(t)$, whose own transform is the \textbf{Dirichlet
kernel}
\[
\Dir_T(f)=\sum_{t\in T} e^{-2\pi i f t}
=\frac{\sin(n\pi f\Delta)}{\sin(\pi f\Delta)}\,e^{-\pi i (n-1) f\Delta}.
\]
Thus $G_n=\Dir_T * G$: observing a finite window \emph{blurs} the true
transform.
\end{definition}

\begin{definition}{Fourier frequencies}{p6-fourierfreq}
The FFT evaluates $G_n$ at the \textbf{Fourier frequencies}
\[
f_k=\frac{k}{\Delta n},\qquad k\in\Ints,\quad -\frac{n}{2}\le k\le \frac{n-1}{2}.
\]
With $\Delta=1$ these are $f_k=k/n$. They are equally spaced over one period
and are the frequencies at which distinct periodogram ordinates are
(approximately) uncorrelated.
\end{definition}

\begin{definition}{Spectral density function (discrete and continuous)}{p6-sdf}
For a discrete-time stationary process $\{X_t\}_{t\in\Ints}$ with ACVS
$\{\acvs_\tau\}\in\ell^1$, the \textbf{spectral density function} is the
discrete-time Fourier transform of the ACVS (with $\Delta=1$):
\[
\sdf(f)=\sum_{\tau\in\Ints}\acvs_\tau\,e^{-2\pi i f\tau},\qquad f\in\Reals .
\]
For a continuous-time process with $\acvs\in L^1$,
$\sdf(f)=\int_{-\infty}^{\infty}\acvs(\tau)e^{-2\pi i f\tau}\,d\tau$. By Fourier
inversion the ACVS is recovered as
$\acvs_\tau=\int_{-1/2}^{1/2}\sdf(f)e^{2\pi i f\tau}\,df$; in particular
$\acvs_0=\Var(X_t)=\int_{-1/2}^{1/2}\sdf(f)\,df$, so the SDF distributes the
total ``power'' $\Var(X_t)$ across frequencies. A valid SDF is real, even
($\sdf(-f)=\sdf(f)$, since $\acvs_{-\tau}=\acvs_\tau$), non-negative, and
periodic with period $1$.
\end{definition}

\begin{definition}{Harmonic (sinusoidal) process}{p6-harmonic}
Let $A>0$ and $\Theta\sim\mathrm{Unif}(-\pi,\pi)$ be independent. The
\textbf{harmonic process} is
\[
X_t=A\cos(\nu t+\Theta),\qquad t\in\Ints,
\]
with $\nu\in\Reals$ the oscillation frequency. Then $\E[X_t]=0$ and
\[
\Cov(X_{t+\tau},X_t)=\tfrac12\,\E[A^2]\cos(\nu\tau),
\]
so it is stationary but its ACVS does \emph{not} decay; in particular
$\{\acvs_\tau\}\notin\ell^1$ and no ordinary SDF exists. Its spectrum is a
pair of lines (jumps) at $\pm\nu/(2\pi)$.
\end{definition}

\begin{definition}{Integrated spectrum (spectral distribution function)}{p6-integspec}
By Herglotz's theorem, every ACVS admits the representation
\[
\acvs_\tau=\int_{-1/2}^{1/2} e^{2\pi i \tau f}\,dS^{(I)}(f),
\]
where $S^{(I)}$, the \textbf{integrated spectrum}, is right-continuous,
bounded, non-decreasing, with $S^{(I)}(-\tfrac12)=0$ and
$S^{(I)}(\tfrac12)=\sigma^2=\acvs_0$. It behaves like a (scaled) cumulative
distribution function over frequency. When $\{\acvs_\tau\}\in\ell^1$ there is a
density $\sdf$ with $S^{(I)}(f)=\int_{-1/2}^{f}\sdf(\lambda)\,d\lambda$, i.e.\
$\sdf=dS^{(I)}/df$, recovering Definition~\ref{def:p6-sdf}.
\end{definition}

\begin{definition}{Periodogram}{p6-periodogram}
For data $X_1,\dots,X_n$ ($T=\{1,\dots,n\}$, $\Delta=1$), the
\textbf{periodogram} is the plug-in estimator obtained by Fourier transforming
the biased sample ACVS:
\[
\hat\sdf^{(p)}(f)=\sum_{\tau\in\Ints}\hat\acvs_\tau\,e^{-2\pi i f\tau},
\qquad
\hat\acvs_\tau=\frac1n\sum_{t=1}^{n-\abs\tau}(X_t-\bar X)(X_{t+\abs\tau}-\bar X)
\]
for $\abs\tau\le n-1$ ($0$ otherwise). Equivalently (Proposition
\ref{prop:p6-periogramJ}) $\hat\sdf^{(p)}(f)=\abs{J(f)}^2$ with
$J(f)=\sqrt{1/n}\sum_{t\in T}(X_t-\bar X)e^{-2\pi i t f}$.
\end{definition}

\begin{definition}{Data taper and tapered periodogram}{p6-taper}
A \textbf{data taper} is a sequence $\{h_t\}_{t\in T}$ with
$\norm{h}_2^2=\sum_t h_t^2=1$. The \textbf{tapered Fourier transform} and
\textbf{tapered periodogram} are
\[
J_h(f)=\sum_{t\in T} h_t\,(X_t-\bar X)\,e^{-2\pi i t f},
\qquad
\hat\sdf^{(p)}_h(f)=\abs{J_h(f)}^2 .
\]
Choosing $h_t=\sqrt{1/n}$ (the constant taper) recovers the ordinary
periodogram. The normalisation $\norm{h}_2^2=1$ keeps the estimator unbiased
for white noise.
\end{definition}

\begin{definition}{Multitaper spectral estimator}{p6-multitaper}
Given $K$ \textbf{orthonormal} tapers $h_k=\{h_{t,k}\}_{t\in T}$ with
$\sum_{t\in T} h_{t,k}h_{t,k'}=\delta_{k,k'}$, the \textbf{multitaper
estimator} averages the individual tapered periodograms,
\[
\hat\sdf^{(mt)}(f)=\frac1K\sum_{k=1}^{K}\hat\sdf^{(p)}_{h_k}(f).
\]
Its expectation is a convolution with the \textbf{average kernel}
$\overline H(f)=\tfrac1K\sum_{k=1}^K\abs{H_k(f)}^2$, and its variance is
$\approx\sdf(f)^2/K$.
\end{definition}

% ======================================================================
\subsection{Theorems \& Propositions}
% ======================================================================

\begin{theorem}{Fourier inversion (continuous and discrete time)}{p6-inversion}
\emph{(Continuous time.)} If $g,G\in L^1$ then the inverse transform recovers
$g$ for almost all $t$:
$g(t)=\int_{-\infty}^{\infty} G(f)\,e^{2\pi i f t}\,df$.
\emph{(Discrete time.)} If $\{g_t\}\in\ell^1$ then $G$ is recovered on one
period as a Fourier series with the domains swapped:
\[
g_t=\int_{-1/(2\Delta)}^{1/(2\Delta)} G(f)\,e^{2\pi i t f}\,df,\qquad
t\in\Delta\Ints .
\]
\textit{Proof idea:} substitute the forward transform into the right-hand side
and use the Fourier orthogonality relation
($\int e^{2\pi i f(t-s)}\,df$ acts as a delta), the integrability hypotheses
justifying the interchange of integration/summation.\\
\textit{Why:} inversion is what makes $(g,G)$ a genuine \emph{Fourier pair} ---
no information is lost. It is the engine behind recovering the ACVS from the SDF
($\acvs_\tau=\int_{-1/2}^{1/2}\sdf(f)e^{2\pi i f\tau}\,df$) and behind every
``time $\leftrightarrow$ frequency'' duality used below.
\end{theorem}

\begin{theorem}{Convolution theorem (continuous time)}{p6-convthm}
For $g,h\in L^1$, if $u=g*h$ then $U=G\cdot H$: the transform of a convolution
is the product of the transforms. If also $G,H\in L^1$ and $v=g\cdot h$, then
dually $V=G*H$.\\[2pt]
\textit{Proof idea:} substitute the convolution integral into the definition of
$U$, swap the order of integration (Fubini, justified by
$\iint\abs{g(x-y)h(y)}\,dx\,dy=\norm g_1\norm h_1<\infty$), and recognise the
shift factor $e^{-2\pi i s f}$ as producing $G(f)$.\\
\textit{Why:} the central computational identity of Fourier analysis ---
``multiplication in one domain is convolution in the other''. It underlies
filtering, the SDF of a filtered process, and the periodogram bias formula.
\end{theorem}

\begin{theorem}{Convolution theorem (discrete time)}{p6-dconvthm}
For $g,h\in\ell^1$ both pairs hold:
\[
h\cdot g \longleftrightarrow H*G,\qquad
h*g \longleftrightarrow H\cdot G .
\]
\textit{Proof idea:} for $u=h*g$, expand $U(f)=\sum_t u_t e^{-2\pi i t f}$,
insert $e^{-2\pi i (t-s)f}e^{-2\pi i s f}$, factor the double sum into
$\big(\sum_s h_s e^{-2\pi i s f}\big)\big(\sum_x g_x e^{-2\pi i x f}\big)$;
interchange is legal because $g,h\in\ell^1$. For $v=h\cdot g$, replace one
factor by its inverse transform and integrate.\\
\textit{Why:} the discrete analogue used to derive the finite-data transform
($G_n=\Dir_T*G$) and the periodogram expectation (a frequency convolution with
the Fej\'er kernel).
\end{theorem}

\begin{theorem}{Aliasing theorem and the Nyquist frequency}{p6-aliasing}
Let $g\in L^1$ be continuous with transform $G$, and let $G$ denote the
transform of the sampled sequence $\{g(t)\}_{t\in\Delta\Ints}$. Under the
summability hypotheses, for every $f$
\[
G(f)=\sum_{k\in\Ints} G(f+k/\Delta).
\]
Sampling \emph{folds} all the continuous-frequency content at
$\{f+k/\Delta\}_{k}$ onto the single frequency $f$. For an SDF this reads
$\sdf(f)=\sum_{k\in\Ints}\sdf(f+k)$ (with $\Delta=1$). The
\textbf{Nyquist frequency} is $1/(2\Delta)$ (i.e.\ $\tfrac12$ when
$\Delta=1$): content above it cannot be recovered.\\[2pt]
\textit{Proof idea:} write $g(t)$ at integer times by both inversion formulas;
split $\int_{-\infty}^{\infty}G(f)e^{2\pi i f t}df$ into period-length blocks,
shift each block by $k/\Delta$ (the phase $e^{2\pi i t k/\Delta}=1$ for
$t\in\Delta\Ints$), swap sum and integral (summability), and match the two
representations of $g(t)$.\\
\textit{Why:} explains why we sample fast enough; if $\sdf$ is negligible
beyond $1/(2\Delta)$ the discrete spectrum faithfully matches the continuous
one, otherwise it is corrupted (poor agreement).
\end{theorem}

\begin{proposition}{FFT complexity}{p6-fft}
The finite Fourier transform at the $n$ Fourier frequencies $\{k/(\Delta n)\}$
can be computed in $O(n\log n)$ time via the \textbf{Fast Fourier Transform},
versus $O(n^2)$ for the naive sum.\\[2pt]
\textit{Proof idea:} divide-and-conquer --- a transform of length $n$ splits
into two of length $n/2$ (even/odd indices), recursing $\log_2 n$ levels with
$O(n)$ work each.\\
\textit{Why:} makes periodogram/spectral computation feasible at scale; one
can also invert the periodogram by FFT to get the sample ACVS cheaply.
\end{proposition}

\begin{theorem}{Herglotz's theorem}{p6-herglotz}
A complex sequence $\{g_t\}_{t\in\Ints}$ is \textbf{non-negative definite} if
and only if there exists a right-continuous, bounded, non-decreasing function
$G^{(I)}$ with $G^{(I)}(-\tfrac12)=0$ such that
\[
g_t=\int_{-1/2}^{1/2} e^{2\pi i t f}\,dG^{(I)}(f),\qquad t\in\Ints .
\]
\textit{Proof idea:} ($\Leftarrow$) any such integral representation yields a
non-negative-definite sequence because $G^{(I)}$ is non-decreasing.
($\Rightarrow$) construct $G^{(I)}$ as the (weak) limit of the non-negative
Ces\`aro/Fej\'er averages of the partial sums $\sum_{\abs t\le N} g_t
e^{-2\pi i t f}$, which form a tight family of measures.\\
\textit{Why:} it is the existence theorem for the spectrum. Since every ACVS is
non-negative definite (a foundational property of any stationary process),
\emph{every} stationary process has an integrated spectrum $S^{(I)}$ --- even
those, like harmonic processes, with no ordinary SDF.
\end{theorem}

\begin{corollary}{Integrated spectrum exists for every stationary process}{p6-integexists}
Applying Herglotz's theorem to the non-negative-definite ACVS gives
$\acvs_\tau=\int_{-1/2}^{1/2}e^{2\pi i \tau f}\,dS^{(I)}(f)$ with
$S^{(I)}(-\tfrac12)=0$, $S^{(I)}(\tfrac12)=\sigma^2$, and $S^{(I)}$
non-decreasing.\\[2pt]
\textit{Why:} guarantees a frequency-domain description always exists; the SDF
is merely the special (absolutely continuous) case $dS^{(I)}=\sdf\,df$.
\end{corollary}

\begin{theorem}{Lebesgue decomposition of the integrated spectrum}{p6-lebesgue}
Any integrated spectrum decomposes uniquely as
\[
S^{(I)}(\cdot)=S^{(I)}_1(\cdot)+S^{(I)}_2(\cdot)+S^{(I)}_3(\cdot),
\]
where each $S^{(I)}_j$ is non-negative, non-decreasing with
$S^{(I)}_j(-\tfrac12)=0$, and
\begin{enumerate}
  \item $S^{(I)}_1$ is \textbf{absolutely continuous}:
        $S^{(I)}_1(f)=\int_{-1/2}^{f}\sdf(\lambda)\,d\lambda$ (the ordinary SDF
        $\sdf$);
  \item $S^{(I)}_2$ is a \textbf{step function} with jumps $\{p_\ell\}$ at
        frequencies $\{f_\ell\}$ of pure sinusoids (a \emph{line spectrum});
  \item $S^{(I)}_3$ is a \textbf{continuous singular} function (pathological,
        of no practical use).
\end{enumerate}
\textit{Proof idea:} the Lebesgue decomposition theorem for monotone functions
/ measures, specialised to the finite interval $[-\tfrac12,\tfrac12]$.\\
\textit{Why:} classifies the possible spectral behaviours: purely continuous
($S^{(I)}_1$ only), purely discrete/line ($S^{(I)}_2$ only), or mixed.
\end{theorem}

\begin{theorem}{Spectral representation theorem}{p6-specrep}
Let $\{X_t\}$ be a real-valued discrete-time stationary process with mean
$\mu$. There exists an \textbf{orthogonal-increment process} $\{Z(f)\}$ on
$[-\tfrac12,\tfrac12]$ with
\[
X_t=\mu+\int_{-1/2}^{1/2} e^{2\pi i f t}\,dZ(f)
\quad\text{(equality in mean square),}
\]
where for $f\ne f'$,
\[
\E[dZ(f)]=0,\qquad
\Var\big(dZ(f)\big)=dS^{(I)}(f),\qquad
\Cov\big(dZ(f),dZ(f')\big)=0 .
\]
\textit{Proof idea:} build $Z$ as the $L^2$ (mean-square) limit of finite
Fourier sums of the $X_t$'s; orthogonality of increments mirrors the
non-decrease of $S^{(I)}$, and the variance of an increment is exactly the mass
$S^{(I)}$ places there.\\
\textit{Why:} the deep structural result: \emph{every} stationary process is a
superposition of uncorrelated random sinusoids $e^{2\pi i f t}\,dZ(f)$, with
``amount of variance'' $dS^{(I)}(f)$ at each frequency. It makes precise the
slogan ``$\sdf(f)\,df$ is the power in $[f,f+df]$''.
\end{theorem}

\begin{proposition}{Periodogram as squared modulus of $J$}{p6-periogramJ}
With $T=\{1,\dots,n\}$ and $Y_t=X_t-\bar X$,
\[
\hat\sdf^{(p)}(f)=\abs{J(f)}^2,\qquad
J(f)=\sqrt{\tfrac1n}\sum_{t\in T} Y_t\,e^{-2\pi i t f}.
\]
\textit{Proof idea:} starting from
$\hat\sdf^{(p)}(f)=\sum_{\tau}\frac1n\sum_t Y_tY_{t+\abs\tau}e^{-2\pi i\tau f}$,
re-index the double lag sum as a free double sum over $s,t\in T$,
\[
\hat\sdf^{(p)}(f)=\frac1n\sum_{s=1}^{n}\sum_{t=1}^{n}
Y_tY_s\,e^{-2\pi i s f}e^{2\pi i t f}
=\Big(\sqrt{\tfrac1n}\sum_t Y_t e^{-2\pi i t f}\Big)
\conj{\Big(\sqrt{\tfrac1n}\sum_s Y_s e^{-2\pi i s f}\Big)}
=\abs{J(f)}^2 .
\]
\textit{Why:} lets the periodogram be computed by a single FFT in
$O(n\log n)$, and makes manifest that $\hat\sdf^{(p)}\ge0$.
\end{proposition}

\begin{theorem}{Expectation of the periodogram: the Fej\'er kernel}{p6-periogexp}
For a zero-mean process,
\[
\E\big[\hat\sdf^{(p)}(f)\big]
=\sum_{\tau\in\Ints} w_\tau\,\acvs_\tau\,e^{-2\pi i\tau f},\qquad
w_\tau=\Big(1-\frac{\abs\tau}{n}\Big)\ind_{[-n,n]}(\tau),
\]
equivalently the frequency convolution
\[
\E\big[\hat\sdf^{(p)}(f)\big]=\int_{-1/2}^{1/2}\sdf(f')\,\Fej_n(f-f')\,df',
\qquad
\Fej_n(f)=
\begin{cases}
\dfrac1n\dfrac{\sin^2(n\pi f)}{\sin^2(\pi f)}, & f\ne0,\\[2.2ex]
n, & f=0,
\end{cases}
\]
the \textbf{Fej\'er kernel}.\\[2pt]
\textit{Proof idea:} $\E[\hat\acvs_\tau]=\big(1-\tfrac{\abs\tau}{n}\big)\acvs_\tau$
for $\abs\tau<n$; transform term by term. The triangular weight
$w_\tau=\tfrac1n(d*\tilde d)_\tau$ is the autocorrelation of the rectangular
window, whose transform is $\abs{\Dir_T(f)}^2/n=\Fej_n(f)$, so multiplication by
$w_\tau$ in time becomes convolution with $\Fej_n$ in frequency.\\
\textit{Why:} quantifies the \textbf{bias}. Because $\Fej_n$ has wide side
lobes that decay slowly, $\E[\hat\sdf^{(p)}]$ is a \emph{blurred} (leaked)
version of $\sdf$; peaks are smeared and power leaks into neighbouring
frequencies. As $n\to\infty$, $\Fej_n$ concentrates at $0$, so the periodogram
is asymptotically unbiased.
\end{theorem}

\begin{theorem}{Variance of the periodogram: inconsistency}{p6-periogvar}
For $f\ne0\ \mathrm{mod}\ \tfrac12$,
\[
\Var\big(\hat\sdf^{(p)}(f)\big)\xrightarrow[n\to\infty]{}\sdf(f)^2 \;>\;0,
\]
and at distinct Fourier frequencies the ordinates are asymptotically
uncorrelated.\\[2pt]
\textit{Proof idea:} for a Gaussian process $J(f)$ is asymptotically complex
Gaussian, so $\hat\sdf^{(p)}(f)=\abs{J(f)}^2$ is asymptotically
$\tfrac12\sdf(f)\chi^2_2=\sdf(f)\cdot\mathrm{Exp}(1)$, whose variance is
$\sdf(f)^2$.\\
\textit{Why:} the periodogram is \textbf{not consistent} --- its variance does
\emph{not} vanish as $n$ grows; more data merely gives more (equally noisy)
ordinates. This is the motivation for smoothing/averaging (multitapering,
or averaging neighbouring ordinates).
\end{theorem}

\begin{theorem}{Expectation of the tapered periodogram}{p6-tapexp}
With a taper $\{h_t\}$ (set $h_t=0$ off $T$) and mean replaced by $\mu$,
\[
\E\big[\hat\sdf^{(p)}_h(f)\big]
=\int_{-1/2}^{1/2}\sdf(f')\,\abs{H(f-f')}^2\,df',
\]
where $H$ is the discrete Fourier transform of $h$.\\[2pt]
\textit{Proof idea:} expand $\hat\sdf^{(p)}_h(f)=\sum_{t,s}h_th_sY_tY_s
e^{-2\pi i f(t-s)}$, take expectations to get $\sum_\tau w_\tau\acvs_\tau
e^{-2\pi i f\tau}$ with $w_\tau=\sum_t h_th_{t+\tau}=(\tilde h*\bar h)_\tau$
(autocorrelation of the taper). By the convolution theorem its transform is
$\tilde H(f)\bar H(f)=\abs{H(f)}^2$, giving the frequency convolution.\\
\textit{Why:} the kernel is now $\abs{H(f)}^2$ instead of the Fej\'er kernel
$\Fej_n$. A well-chosen taper (e.g.\ a dpss / Slepian taper) makes
$\abs{H(f)}^2$ far more concentrated near $0$ with much lower side lobes, so
\textbf{tapering reduces bias / leakage}.
\end{theorem}

\begin{corollary}{Unbiasedness of the tapered periodogram for white noise}{p6-tapunbias}
If $\{X_t\}$ is white noise with variance $\sigma^2$ (so $\sdf\equiv\sigma^2$)
and $\norm{h}_2^2=1$, then $\E[\hat\sdf^{(p)}_h(f)]=\sigma^2=\sdf(f)$ for all
$n$ and all $f$.\\[2pt]
\textit{Proof idea:} the kernel integrates out,
$\E[\hat\sdf^{(p)}_h(f)]=\sigma^2\int_{-1/2}^{1/2}\abs{H(f')}^2\,df'
=\sigma^2\sum_{t}h_t^2=\sigma^2\norm{h}_2^2$ (Parseval), which equals
$\sigma^2$ exactly when $\norm h_2^2=1$.\\
\textit{Why:} pins down the normalisation $\norm h_2^2=1$ in the definition of
a taper, and shows tapering costs nothing in bias for flat spectra.
\end{corollary}

\begin{theorem}{Variance reduction by multitapering}{p6-mtvar}
For $K$ orthonormal tapers, asymptotically
\[
\Cov\big(\hat\sdf^{(p)}_{h_k}(f),\hat\sdf^{(p)}_{h_{k'}}(f)\big)
\to\sdf(f)^2\,\delta_{k,k'},
\quad\text{hence}\quad
\Var\big(\hat\sdf^{(mt)}(f)\big)\approx\frac{\sdf(f)^2}{K}.
\]
\textit{Proof idea:} orthogonality of the tapers makes the $K$ individual
tapered periodograms asymptotically uncorrelated, each with variance
$\sdf(f)^2$; averaging $K$ uncorrelated terms divides the variance by $K$.\\
\textit{Why:} \textbf{multitapering restores consistency}: as $n\to\infty$ one
may let $K\to\infty$ (slowly) so $\Var\to0$. It is the standard fix for the
periodogram's non-vanishing variance, with only a modest bias increase.
\end{theorem}

% ======================================================================
% ======================================================================

% ======================================================================
% ======================================================================

% ---------------------- NEW EXERCISES ----------------------------------

